\section{Related Work}

\textbf{Detection on these benchmarks.} Guo et al.~\cite{Guo2023close} introduce HC3 and report a RoBERTa detector at 99.82 F1 on English full text, establishing both the corpus and the accuracy regime the field has worked in since. Su et al.~\cite{Su2023plus} show that regime is narrower than it appears, since moving from question answering to semantic-invariant tasks drops accuracy to roughly 91.7\%, though their evaluation stays within ChatGPT-generated text. Zhou~\cite{Zhou2024exploiting} reaches competitive accuracy on DAIGT V2 with classical models over character n-grams and no transformer at all, which is consistent with our DAIGT V2 finding but is reported without a surface-content separation that would explain it. Kubrusly et al.~\cite{Kubrusly2025detecting} likewise report strong bag-of-words results on HC3, and likewise do not isolate what the features encode. Ansary~\cite{Ansary2026multirepresentation} and Lamsiyah et al.~\cite{Lamsiyah2025daigt} report near-ceiling results on DAIGT-family data, in both cases with in-distribution evaluation only. Annepaka et al.~\cite{Annepaka2026synergizing} combine explicit linguistic features with a transformer and reach 99.45\% on HC3, a design that mixes the two information sources our decomposition separates.

\textbf{Documented artefacts.} Tian et al.~\cite{Tian2023multiscale} identify the HC3 whitespace convention, observe that a detector could score well by exploiting it, and release a cleaning kit. Their contribution is the artefact and the remedy. They do not quantify how much of the corpus's total separability surface form carries, and they do not test whether a given detector can represent the cue at all, which is the question our tokenisation control answers and which turns out to change the conclusion for BERT. Baidya et al.~\cite{Baidya2026detecting} document a length confound in HC3 and control it by length-matching, again treating one artefact in isolation. Ardeshirifar~\cite{Ardeshirifar2025comparing} standardises punctuation and contractions before comparing handcrafted features to deep models, establishing the practice as methodologically necessary without measuring what it removes. Park et al.~\cite{Park2024investigating} show HC3-trained detectors rely on prompt-specific collection shortcuts and construct instructions exploiting them. Our decomposition is complementary to all four. Where they identify individual cues, we measure the aggregate, which is the quantity a practitioner comparing two benchmarks actually needs.

\textbf{Robustness and its controls.} Antoun et al.~\cite{Antoun2023towards} attack HC3-trained detectors with misspellings and homoglyphs and report a fall from 99.88 F1 to 33.57\% accuracy on a hand-built adversarial set. Huang et al.~\cite{Huang2024generated} report comparable degradation under systematic character perturbation. Neither reports a label-free control, and our results suggest such collapses cannot be separated from a generic degenerate response to unreadable input without one, since our own models label random character strings as human with near-maximal confidence. This is a methodological gap rather than a dispute about their measurements. Lai et al.~\cite{Lai2024adaptive} show that ensembling fine-tuned detectors preserves in-distribution accuracy while still degrading under distribution shift, so the fragility is not an artefact of single-model variance. Zhang et al.~\cite{Zhang2024detecting} step outside per-document classification altogether, framing detection as a two-sample test, which sidesteps the single-document shortcuts studied here at the cost of needing a population. Borile and Abrate~\cite{Borile2025generalize} take the complementary mechanistic route, ablating roughly twenty feed-forward neurons to improve out-of-distribution accuracy by up to 6.9 points, evidence that some in-distribution performance rests on localised corpus-specific features.

\textbf{Cross-corpus evaluation.} Alikhanov et al.~\cite{Alikhanov2026generated} merge HC3 and DAIGT V2 into a single 124,195-sample corpus and evaluate under a topic-held-out split, reporting 82.87\% to 88.86\% accuracy against the near-ceiling numbers usual in-distribution. Their design asks whether one model can serve both corpora. Ours asks what each corpus is separable by, which is why the two results sit alongside rather than replacing one another. Lu et al.~\cite{Lu2023large} show detectors can be evaded by prompting alone, a threat model orthogonal to the corpus-property question studied here.
